\documentclass[12pt, a4paper]{article}

\usepackage{amsmath, amssymb, amsfonts}
\usepackage{float}
\usepackage{graphicx}
\usepackage{listings}
\usepackage{xcolor}
\usepackage[linesnumbered, ruled]{algorithm2e}
\usepackage{subcaption}
\usepackage[hidelinks]{hyperref}
% Each algorithm is written twice: the algorithm2e block (what the PDF typesets)
% and a plain-text Verbatim block right after it (what the web shows — pandoc
% renders Verbatim as a copyable <pre> code block). \excludecomment makes pdflatex
% SKIP the Verbatim blocks, so the PDF shows only the algorithm2e version; pandoc
% conversely drops each algorithm2e to an (empty, CSS-hidden) div and keeps the
% Verbatim. Because pdflatex never sees the Verbatim, its Unicode maths glyphs
% (⟨⟩ ← τ ∈ ∞ …) need no inputenc/literate setup.
\usepackage{comment}
\excludecomment{Verbatim}

\definecolor{codegreen}{rgb}{0,0.6,0}
\definecolor{codegray}{rgb}{0.5,0.5,0.5}
\definecolor{codepurple}{rgb}{0.58,0,0.82}
\hypersetup{linktoc=all}

\lstset{
    language=Python,
    basicstyle=\ttfamily\small,
    keywordstyle=\color{blue},
    commentstyle=\color{codegreen},
    stringstyle=\color{codepurple},
    breaklines=true,
    numbers=left,
    numberstyle=\tiny\color{codegray},
    frame=single,
    showstringspaces=false
}

\title{AI Masterclass Test 1}
\author{}
\date{}


\begin{document}

\section{A* Search}
Minimises both path cost so far \textit{and} estimated cost to goal,
\[f(s) = g(s) + h(s)\]

\paragraph{Admissible heuristics} A heuristic function is admissible (optimistic) if
it never overestimates the true cost to the goal.
\[0\leq h(s) \leq h^*(s), \quad \forall s\]
where $h^*(s)$ is the true lowest cost to from $s$ to the goal.

\paragraph{Theorem} If $h(s)$ is admissible, A* is guaranteed to find an optimal path to the goal.\newline 
Because A* never overestimates, A* never ignores a shorter path in favour of a longer one. 
So when A* reaches the goal for the first time, it is guaranteed to be the optimal one.


\paragraph{Example} 2D grid: Manhatten distance (for strictly vertical and horizontal movements)
\[h(s) = |x_s - x_g| + |y_s - y_g|\]
where $(x_s, y_s)$ is coordinate of state $s$ and $(x_g, y_g)$ is coordinate of goal state.\newline
The actual steps needed is at least equal to the number of horizontal steps + vertical steps.

\paragraph{Example} Unrestricted movement in 2D space: Euclidean distance
\[h(s) = \sqrt{(x_s - x_g)^2 + (y_s - y_g)^2}\]

\paragraph{Example} 8-puzzle
\begin{figure}[H]
    \centering
    \includegraphics[width=0.5\textwidth]{images/8-puzzle.png} 
    \caption{$h(s)$ = number of misplaced tiles or sum of Manhatten distances of each tile to their goal position}
    \label{fig:my_handwritten_notes}
\end{figure}

\end{document}
